\textbf{Phần 1. Trắc nghiệm lựa chọn}
\setcounterpageref{ex}{0}

\begin{ex} %[3X15N74]
Khi tung một đồng xu gồm mặt sấp và mặt ngửa, sự kiện nào có thể xảy ra?
\choice
{Chắc chắn xuất hiện mặt sấp}
{Không thể xuất hiện mặt sấp}
{Chắc chắn xuất hiện cả 2 mặt cùng lúc}
{\True Có thể xuất hiện mặt sấp hoặc mặt ngửa}
\loigiai{
Khi tung đồng xu thì có thể xuất hiện mặt sấp hoặc mặt ngửa.\\
\textbf{Đáp án của câu hỏi là D}
}
\end{ex}

\begin{ex} %[3X15N74]
Trong hộp chỉ có toàn bóng màu đỏ. Lấy ngẫu nhiên 1 quả bóng ra khỏi hộp thì sự kiện nào là đúng?
\choice
{\True Chắc chắn lấy được bóng màu đỏ}
{Có thể lấy được bóng màu xanh}
{Không thể lấy được bóng màu đỏ}
{Chắc chắn lấy được bóng màu vàng}
\loigiai{
Vì trong hộp chỉ có bóng đỏ nên chắc chắn lấy được bóng đỏ.\\
\textbf{Đáp án của câu hỏi là A}
}
\end{ex}

\begin{ex} %[3X15H74]
Trong hộp chỉ có bóng xanh và bóng vàng. Lấy ngẫu nhiên 1 quả bóng thì sự kiện \textit{"Lấy được bóng màu đỏ"} là
\choice
{Chắc chắn xảy ra}
{\True Không thể xảy ra}
{Có thể xảy ra}
{Luôn luôn xảy ra}
\loigiai{
Trong hộp không có bóng đỏ nên sự kiện không thể xảy ra.\\
\textbf{Đáp án của câu hỏi là B}
}
\end{ex}

\begin{ex} %[3X15H74]
Gieo một con xúc xắc 6 mặt (từ 1 đến 6 chấm). Sự kiện \textit{"Số chấm xuất hiện là 7"} là sự kiện
\choice
{Chắc chắn xảy ra}
{Có thể xảy ra}
{Rất dễ xảy ra}
{\True Không thể xảy ra}
\loigiai{
Xúc xắc chỉ có từ 1 đến 6 chấm nên không thể xuất hiện 7 chấm.\\
\textbf{Đáp án của câu hỏi là D}
}
\end{ex}

\begin{ex} %[3X15H74]
Gieo một con xúc xắc 6 mặt. Sự kiện \textit{"Số chấm xuất hiện là số chẵn"} là sự kiện
\choice
{Chắc chắn xảy ra}
{\True Có thể xảy ra}
{Không thể xảy ra}
{Luôn luôn đúng}
\loigiai{
Có thể xuất hiện các số chẵn (2, 4, 6) hoặc số lẻ (1, 3, 5).\\
\textbf{Đáp án của câu hỏi là B}
}
\end{ex}

\begin{ex} %[3X15H74]
Trong túi có 3 viên bi xanh và 2 viên bi đỏ. Lấy ra cùng lúc 3 viên bi. Sự kiện nào sau đây \textbf{chắc chắn} xảy ra?
\choice
{Lấy được 3 viên bi đỏ}
{Lấy được 3 viên bi vàng}
{\True Lấy được ít nhất 1 viên bi xanh}
{Lấy được 2 viên bi vàng}
\loigiai{
Vì chỉ có 2 bi đỏ nên khi lấy 3 viên chắc chắn phải có ít nhất 1 bi xanh.\\
\textbf{Đáp án của câu hỏi là C}
}
\end{ex}

\begin{ex} %[3X15V74]
Quay vòng quay may mắn chia đều 4 phần: Đỏ, Xanh, Vàng, Trắng. Mũi tên dừng lại ở ô màu Tím là sự kiện
\choice
{\True Không thể xảy ra}
{Chắc chắn xảy ra}
{Có thể xảy ra}
{Rất có thể xảy ra}
\loigiai{
Vòng quay không có màu Tím nên không thể xảy ra.\\
\textbf{Đáp án của câu hỏi là A}
}
\end{ex}

\begin{ex} %[3X15V74]
Một hộp có 10 chiếc thẻ ghi số từ 1 đến 10. Rút ngẫu nhiên 1 thẻ. Sự kiện \textit{"Rút được số bé hơn 11"} là sự kiện
\choice
{Có thể xảy ra}
{\True Chắc chắn xảy ra}
{Không thể xảy ra}
{Hiếm khi xảy ra}
\loigiai{
Các số từ 1 đến 10 đều bé hơn 11 nên chắc chắn xảy ra.\\
\textbf{Đáp án của câu hỏi là B}
}
\end{ex}

\begin{ex} %[3X15V74]
Rút ngẫu nhiên 1 thẻ trong các thẻ ghi số $2, 4, 6, 8$. Sự kiện \textit{"Rút được thẻ ghi số lẻ"} là sự kiện
\choice
{Chắc chắn xảy ra}
{Có thể xảy ra}
{\True Không thể xảy ra}
{Tất nhiên xảy ra}
\loigiai{
Tất cả các thẻ đều là số chẵn nên không thể rút được số lẻ.\\
\textbf{Đáp án của câu hỏi là C}
}
\end{ex}

\begin{ex} %[3X15V74]
Ngày mai Mặt Trời mọc ở hướng Đông là sự kiện
\choice
{\True Chắc chắn xảy ra}
{Có thể xảy ra}
{Không thể xảy ra}
{Chưa biết được}
\loigiai{
Mặt Trời luôn mọc ở hướng Đông nên sự kiện này chắc chắn xảy ra.\\
\textbf{Đáp án của câu hỏi là A}
}
\end{ex}

\begin{ex} %[3X15V74]
Hôm nay là thứ Hai, ngày mai là thứ Tư là sự kiện
\choice
{Chắc chắn xảy ra}
{Có thể xảy ra}
{Hiếm khi xảy ra}
{\True Không thể xảy ra}
\loigiai{
Sau thứ Hai chắc chắn là thứ Ba (không thể là thứ Tư).\\
\textbf{Đáp án của câu hỏi là D}
}
\end{ex}

\begin{ex} %[3X15C74]
Trong hộp có 5 viên kẹo dâu và 5 viên kẹo cam. Bạn Bình muốn lấy chắc chắn có ít nhất 1 viên kẹo dâu thì Bình cần lấy ra ít nhất bao nhiêu viên kẹo (không nhìn vào hộp)?
\choice
{$5$ viên}
{$7$ viên}
{\True $6$ viên}
{$10$ viên}
\loigiai{
Trường hợp xấu nhất lấy phải 5 viên kẹo cam, lấy thêm 1 viên nữa chắc chắn là kẹo dâu: $5 + 1 = 6$ viên.\\
\textbf{Đáp án của câu hỏi là C}
}
\end{ex}

\textbf{Phần 2. Trắc nghiệm đúng sai}
\setcounterpageref{ex}{0}

\begin{ex} %[3X15H74]
Xét tính đúng sai về khả năng xảy ra của các sự kiện:
\choiceTF
{\True Gieo xúc xắc 6 mặt, số chấm xuất hiện có thể là 5.}
{\True Gieo xúc xắc 6 mặt, số chấm xuất hiện chắc chắn bé hơn 7.}
{Gieo xúc xắc 6 mặt, số chấm xuất hiện chắc chắn là 6.}
{\True Số chấm xuất hiện có thể là 6 (hoặc 1, 2, 3, 4, 5), không phải chắc chắn.}
\loigiai{
\begin{itemize}
    \item Có thể là 5: Đúng.
    \item Chắc chắn bé hơn 7 (từ 1 đến 6): Đúng.
    \item Không chắc chắn là 6. Ý 3 Sai, Ý 4 Đúng.
\end{itemize}
\begin{itemize}
    \item a) \textbf{Đúng}.
    \item b) \textbf{Đúng}.
    \item c) \textbf{Sai}.
    \item d) \textbf{Đúng}.
\end{itemize}
}
\end{ex}

\begin{ex} %[3X15V74]
Trong túi có 4 quả bóng xanh và 3 quả bóng đỏ. Xét tính đúng sai khi lấy ngẫu nhiên 1 quả bóng:
\choiceTF
{\True Có thể lấy được quả bóng màu xanh.}
{\True Có thể lấy được quả bóng màu đỏ.}
{Chắc chắn lấy được quả bóng màu xanh.}
{\True Không thể lấy được quả bóng màu vàng.}
\loigiai{
\begin{itemize}
    \item Có thể xanh: Đúng.
    \item Có thể đỏ: Đúng.
    \item Có thể là bi đỏ nên không chắc chắn lấy được bi xanh. Sai.
    \item Không có bóng vàng nên không thể lấy được bóng vàng: Đúng.
\end{itemize}
\begin{itemize}
    \item a) \textbf{Đúng}.
    \item b) \textbf{Đúng}.
    \item c) \textbf{Sai}.
    \item d) \textbf{Đúng}.
\end{itemize}
}
\end{ex}

\textbf{Phần 3. Trả lời ngắn}
\setcounterpageref{ex}{0}

\begin{ex} %[3X15H74]
Gieo một con xúc xắc 6 mặt, có bao nhiêu khả năng có thể xuất hiện về số chấm (từ 1 đến 6)?
\shortans{6}
\loigiai{
\textbf{Đáp số:} $6$\\
6 khả năng (1, 2, 3, 4, 5, 6 chấm).
}
\end{ex}

\begin{ex} %[3X15H74]
Khi tung một đồng xu, có bao nhiêu khả năng có thể xuất hiện về mặt của đồng xu?
\shortans{2}
\loigiai{
\textbf{Đáp số:} $2$\\
2 khả năng (mặt sấp hoặc mặt ngửa).
}
\end{ex}

\begin{ex} %[3X15V74]
Trong hộp có 4 bút bi xanh, 3 bút bi đỏ và 2 bút bi đen. Có tất cả bao nhiêu loại màu bút có thể lấy ra?
\shortans{3}
\loigiai{
\textbf{Đáp số:} $3$\\
3 màu (xanh, đỏ, đen).
}
\end{ex}

\begin{ex} %[3X15V74]
Gieo một con xúc xắc 6 mặt, có bao nhiêu khả năng xuất hiện mặt có số chấm là số chẵn?
\shortans{3}
\loigiai{
\textbf{Đáp số:} $3$\\
3 khả năng (2, 4, 6 chấm).
}
\end{ex}

\begin{ex} %[3X15V74]
Một hộp có 10 viên bi đỏ và 1 viên bi xanh. Cần lấy ra ít nhất bao nhiêu viên bi để chắc chắn có ít nhất 1 viên bi màu đỏ?
\shortans{2}
\loigiai{
\textbf{Đáp số:} $2$\\
Lấy $1 + 1 = 2$ viên bi.
}
\end{ex}

\begin{ex} %[3X15C74]
Trong túi có 6 viên kẹo cam và 4 viên kẹo dâu. Cần lấy ra ít nhất bao nhiêu viên kẹo để chắc chắn có ít nhất 1 viên kẹo dâu?
\shortans{7}
\loigiai{
\textbf{Đáp số:} $7$\\
$6 + 1 = 7$ viên kẹo.
}
\end{ex}
